% ---------------------------------------------------------------
%   dvilualatex -interaction=nonstopmode header-molecule.tex
%   dvisvgm --no-fonts --exact-bbox -o header-molecule.raw.svg header-molecule.dvi
%   python3 ../scripts/animate_svg.py header-molecule.raw.svg -o header-molecule-light.svg
% ---------------------------------------------------------------
\documentclass[crop,border=6pt]{standalone}
\def\pgfsysdriver{pgfsys-dvisvgm.def}
\usepackage{chemfig}

\begin{document}

\setchemfig{
  atom sep        = 2.6em,   % 結合長
  bond style      = {line width=1.15pt},
  double bond sep = 3.4pt,
}

% ---- 描く分子はこの1行だけ差し替えればよい ----------------------
% トリプタミン
% 縮環は「小さい環をブランチ (*5(...)) として書く」記法でないと閉じない。
% また *6(-=-(...)) のように2本引いた後の頂点に付けるのが安定。
\chemfig{*6(-=-(*5(-NH-=(-[:60]-[:0]-[:60]NH_2)-))=-=)}

% 他の例:
% ベンゼン        \chemfig{*6(-=-=-=)}
% サリチル酸      \chemfig{*6(=-=(-C(=[:30]O)-[:-30]OH)-(-OH)=-)}
% カフェイン      \chemfig{*6((=O)-N(-)-(*5(-N=-N(-)-))=-(=O)-N(-)-)}
% -----------------------------------------------------------------

\end{document}
